% ===========================================================================
%  MANUEL COMPLET : JavaScript - PHP - Java
%  Document d'auto-apprentissage, avec pratique, outils et liens.
%
%  Compilation (une seule commande suffit) :
%     pdflatex main.tex
%     pdflatex main.tex        <- 2e passage : table des matieres + ref croisees
%  ou, mieux :
%     latexmk -pdf -interaction=nonstopmode main.tex
%  ou, sans TeX Live installe :
%     tectonic main.tex
%
%  Si un paquet manque (erreur "File `xxx.sty' not found"), installe-le :
%     sudo apt install texlive-latex-extra texlive-fonts-extra texlive-lang-french
%     tlmgr install libertine inconsolata cleveref titletoc menu
% ===========================================================================

\documentclass[12pt,a4paper,twoside,openright]{book}

% ------------------------------ Encodage / langue --------------------------
\usepackage[T1]{fontenc}
\usepackage[utf8]{inputenc}
\usepackage[french]{babel}

% ------------------------------ Mise en page ------------------------------
\usepackage[left=2.1cm,right=2.1cm,top=2.3cm,bottom=2.4cm,headheight=15pt]{geometry}
\usepackage{microtype}

% ------------------------------ Polices -----------------------------------
\usepackage{libertine}
\usepackage{amsmath}
\usepackage{amssymb}
\usepackage[scaled=0.92]{inconsolata}

% ------------------------------ Liens / navigation ------------------------
\usepackage[table]{xcolor}
\usepackage{hyperref}
\definecolor{linkblue}{HTML}{0B4F8A}
\definecolor{pathgreen}{HTML}{1F6B3A}
\hypersetup{
  colorlinks=true, linkcolor=linkblue, filecolor=linkblue, urlcolor=linkblue,
  citecolor=linkblue, bookmarksnumbered=true, bookmarksopen=true, unicode=true,
  pdftitle={Manuel complet : JavaScript, PHP et Java},
  pdfauthor={Manuel d'auto-apprentissage},
  pdfsubject={Programmation : JavaScript, PHP, Java},
  pdfkeywords={JavaScript, Node.js, PHP, MySQL, Java, Spring Boot, tutoriel}
}
\usepackage[capitalise,nameinlink]{cleveref}

% ------------------------------ Tableurs / listes -------------------------
\usepackage{graphicx}
\usepackage{booktabs}
\usepackage{longtable}
\usepackage{array}
% X = colonne a largeur fixe (tabularx/ltablex sont abandonnes :
% leur re-mesure du corps cassait \verb a l'interieur des cellules)
\newcolumntype{X}{>{\raggedright\arraybackslash}p{0.30\textwidth}}
% tabularx n'est plus charge : on le remappe proprement sur longtable
\newenvironment{tabularx}[2][\textwidth]{\rowcolors{2}{white}{tralight}\small
  \setlength{\tabcolsep}{5pt}\begin{longtable}{#2}}
  {\end{longtable}\normalfont\setlength{\tabcolsep}{6pt}}
\usepackage{enumitem}
\setlist[itemize]{topsep=3pt,itemsep=2pt,parsep=0pt,leftmargin=1.5em}
\setlist[enumerate]{topsep=3pt,itemsep=2pt,parsep=0pt,leftmargin=1.6em}
\setlist[description]{topsep=3pt,itemsep=3pt,parsep=0pt,font=\bfseries}
\usepackage{multicol}
\usepackage{menukeys}
\usepackage{titletoc}
\usepackage{makeidx}
\makeindex
\titlecontents{section}[3.2em]{\addvspace{2pt}\small}
  {\contentslabel{2.0em}}{}{\titlerule*[0.75pc]{.}\contentspage}

% ------------------------------ Couleurs du theme -------------------------
\definecolor{codebg}{HTML}{F6F8FA}
\definecolor{codeframe}{HTML}{D0D7DE}
\definecolor{ctermbg}{HTML}{1B2430}
\definecolor{ctermfg}{HTML}{E6EDF3}
\definecolor{codekw}{HTML}{0550AE}
\definecolor{codecom}{HTML}{5E7C40}
\definecolor{codestr}{HTML}{A31515}
\definecolor{themecolor}{HTML}{0B4F8A}
\definecolor{jsyellow}{HTML}{F0DB4F}
\definecolor{phppurple}{HTML}{777BB4}
\definecolor{javared}{HTML}{E76F00}
\definecolor{boxtitle}{HTML}{0B4F8A}
\definecolor{objbg}{HTML}{EAF3FB}\definecolor{objfr}{HTML}{8FC1E8}
\definecolor{notebg}{HTML}{EFF6EE}\definecolor{notefr}{HTML}{A7C9A0}
\definecolor{tipbg}{HTML}{FFF6DE}\definecolor{tipfr}{HTML}{E3C06A}
\definecolor{warnbg}{HTML}{FDECEC}\definecolor{warnfr}{HTML}{E6A0A0}
\definecolor{exobg}{HTML}{F1EEFB}\definecolor{exofr}{HTML}{B7A6E0}
\definecolor{linksbg}{HTML}{EEF3F7}\definecolor{linksfr}{HTML}{9BB8CE}
\definecolor{theadbg}{HTML}{0B4F8A}
\definecolor{tralight}{HTML}{F1F5F9}

% ------------------------------ Code (listings) ---------------------------
\usepackage{listings}
\lstdefinestyle{base}{
  basicstyle=\ttfamily\small,
  columns=fullflexible,
  keepspaces=true,
  breaklines=true,
  breakatwhitespace=true,
  showstringspaces=false,
  tabsize=2,
  numbers=none,
  frame=single,
  rulecolor=\color{codeframe},
  framerule=0.4pt,
  framesep=4pt,
  backgroundcolor=\color{codebg},
  xleftmargin=6pt,
  xrightmargin=6pt,
  aboveskip=8pt,
  belowskip=6pt,
  commentstyle=\color{codecom},
  keywordstyle=\color{codekw}\bfseries,
  stringstyle=\color{codestr},
  identifierstyle=\color{black},
  captionpos=b
}
\lstdefinestyle{jscode}{style=base,language=JavaScript,morekeywords={async,await,of,as,static,let,const,yield,typeof,instanceof,void,delete,in,export,import,from,catch,finally,throw,undefined,NaN,Infinity,extends,super,get,set,Symbol,BigInt,arguments,debugger}}
\lstdefinestyle{phpcode}{style=base,language=[5.6]PHP,morekeywords={function,class,interface,trait,enum,abstract,final,extends,implements,namespace,use,const,public,protected,private,static,readonly,match,fn,global,insteadof,as,print,unset,isset,echo,require,require_once,include,include_once,declare,self,parent,__DIR__,__FILE__,__LINE__}}
\lstdefinestyle{javacode}{style=base,language=Java,morekeywords={var,record,sealed,permits,yield,instanceof,enum,interface,class,abstract,final,static,public,protected,private,synchronized,volatile,transient,native,implements,extends,package,import,throws,throw,try,catch,finally,super,this,assert,switch,boolean,char,byte,short,int,long,float,double,void}}
\lstdefinestyle{confcode}{style=base,morecomment=[l]{//},morecomment=[l]{;}}
\lstdefinestyle{sqlcode}{style=base,language=SQL,morekeywords={SELECT,INSERT,UPDATE,DELETE,CREATE,ALTER,DROP,TABLE,INDEX,VIEW,TRIGGER,PRIMARY,KEY,FOREIGN,REFERENCES,UNIQUE,NOT,NULL,DEFAULT,AUTO_INCREMENT,ENGINE,VARCHAR,INT,BIGINT,TEXT,BLOB,DATE,DATETIME,TIMESTAMP,DECIMAL,BOOLEAN,JOIN,INNER,LEFT,RIGHT,OUTER,ON,AS,WHERE,GROUP,ORDER,BY,HAVING,LIMIT,OFFSET,DISTINCT,AND,OR,IN,LIKE,BETWEEN,IS,COUNT,SUM,AVG,MIN,MAX,IF,EXISTS,UNION,ALL,VALUES,SET,TRANSACTION,COMMIT,ROLLBACK,GRANT,DESC,ASC,CASCADE,CONSTRAINT,CHECK,SHOW,USE,DATABASE}}
\lstdefinestyle{shellcode}{style=base,language=bash}
\lstdefinestyle{jsoncode}{style=base,morecomment=[l]{//}}
\lstdefinestyle{ymlcode}{style=base}

\lstset{style=base}

% ------------------------------ Boites (tcolorbox) -------------------------
\usepackage[most]{tcolorbox}

% (l'ancienne cle "codestyle" avec "listing only" est supprimee : c'est le
% mecanisme de fichier temporaire output.listing qui faisait planter la
% compilation, et il n'est plus necessaire)

% Chaque bloc de code = un environnement listings autonome (il casse les
% lignes tout seul, se coupe proprement en fin de page, et ne cree aucun
% fichier temporaire).
\lstnewenvironment{codejs}{\lstset{style=jscode}}{}
\lstnewenvironment{codephp}{\lstset{style=phpcode}}{}
\lstnewenvironment{codejava}{\lstset{style=javacode}}{}
\lstnewenvironment{codesql}{\lstset{style=sqlcode}}{}
\lstnewenvironment{codeconf}{\lstset{style=confcode}}{}
\lstnewenvironment{codejson}{\lstset{style=jsoncode}}{}
\lstnewenvironment{codeyml}{\lstset{style=ymlcode}}{}
\lstnewenvironment{codeplain}{\lstset{style=base}}{}
\lstnewenvironment{terminal}{\lstset{style=shellcode,frame=single,
  rulecolor=\color{ctermbg},backgroundcolor=\color{ctermbg},
  basicstyle=\ttfamily\small\color{ctermfg},
  keywordstyle=\color{ctermfg}\bfseries,commentstyle=\color{ctermfg}}}{}

\newtcbox{\srcbox}{on line, colback=tralight, colframe=tralight, boxrule=0pt,
  arc=0.8mm, left=1.2mm,right=1.2mm,top=0.5mm,bottom=0.7mm, size=tight,
  fontupper=\scriptsize}

\newtcolorbox{objectif}[1][]{enhanced, breakable, colback=objbg, colframe=objfr,
  boxrule=0.5pt, arc=1.5mm, left=3mm,right=3mm,top=2mm,bottom=2mm,
  fonttitle=\bfseries, coltitle=white, colbacktitle=boxtitle,
  title={Objectifs du chapitre},#1}

\newtcolorbox{note}{enhanced, breakable, colback=notebg, colframe=notefr,
  boxrule=0.5pt, arc=1.5mm, left=3mm,right=3mm,top=2mm,bottom=2mm,
  before upper={\small\textbf{Note.}\enspace}}

\newtcolorbox{astuce}{enhanced, breakable, colback=tipbg, colframe=tipfr,
  boxrule=0.5pt, arc=1.5mm, left=3mm,right=3mm,top=2mm,bottom=2mm,
  before upper={\small\textbf{Astuce de pro.}\enspace}}

\newtcolorbox{piege}{enhanced, breakable, colback=warnbg, colframe=warnfr,
  boxrule=0.5pt, arc=1.5mm, left=3mm,right=3mm,top=2mm,bottom=2mm,
  before upper={\small\textbf{Pi\`ege \`a \`eviter.}\enspace}}

\newtcolorbox{exercice}[1][]{enhanced, breakable, colback=exobg, colframe=exofr,
  boxrule=0.5pt, arc=1.5mm, left=3mm,right=3mm,top=2mm,bottom=2mm,
  fonttitle=\bfseries, coltitle=white, colbacktitle=boxtitle,
  title={Exercices guid\'es},#1}

\newtcolorbox{ressources}[1][]{enhanced, breakable, colback=linksbg, colframe=linksfr,
  boxrule=0.5pt, arc=1.5mm, left=3mm,right=3mm,top=2mm,bottom=2mm,
  fonttitle=\bfseries, coltitle=white, colbacktitle=boxtitle,
  title={Ressources et outils cit\'es},before upper=\small,#1}

\newtcolorbox{projet}[1][]{enhanced, breakable, colback=tipbg, colframe=tipfr,
  boxrule=0.5pt, arc=1.5mm, left=3mm,right=3mm,top=2mm,bottom=2mm,
  fonttitle=\bfseries, coltitle=white, colbacktitle=javared,
  title={Projet \`a r\'ealiser},#1}

% ------------------------------ Macros pratiques ---------------------------
\newcommand{\fichier}[1]{\texttt{\color{pathgreen}#1}}
\newcommand{\rep}[1]{\texttt{\color{pathgreen}#1/}}
\providecommand{\cmd}[1]{}%
\renewcommand{\cmd}[1]{\srcbox{\texttt{\detokenize{#1}}}}
\newcommand{\link}[2]{\href{#2}{#1}}
\newcommand{\src}[1]{\href{#1}{\textcolor{linkblue}{\textsuperscript{[src]}}}}
\newcommand{\langjs}{\textcolor{javared}{\textbf{JavaScript}}}
\newcommand{\langphp}{\textcolor{phppurple}{\textbf{PHP}}}
\newcommand{\langjava}{\textcolor{themecolor}{\textbf{Java}}}
\newcommand{\etape}[1]{\par\medskip\noindent\textbf{\large #1}\par\smallskip}
\DeclareRobustCommand{\clic}[1]{\textcolor{linkblue}{\textsf{#1}}}

% tableaux zebras : tabz = longtable (largeurs fixes) + lignes alternees
\newenvironment{tabz}[1]{\rowcolors{2}{white}{tralight}%
  \small\setlength{\tabcolsep}{5pt}\setlength{\LTpre}{3pt}\setlength{\LTpost}{3pt}%
  \begin{longtable}{@{}#1@{}}\toprule}
  {\bottomrule\end{longtable}\normalfont\setlength{\tabcolsep}{6pt}}
\newcommand{\thd}[1]{\textbf{\color{white}#1}}
\newcommand{\headrow}{\rowcolor{theadbg}}
\makeatletter
\newenvironment{tablez}{\@ifnextchar[{\@tablezwith}{\@tablezwith[XXX]}}{\end{tabz}}
\newcommand{\@tablezwith}[1]{\begin{tabz}{#1}}
\makeatother

% ------------------------------ Doc ---------------------------------------
\begin{document}
\mainmatter
% ============================================================================
\chapter{M\'ethode d'apprentissage et plan de 12 semaines}
\label{ch-methode}
% ============================================================================

Savoir quoi apprendre ne suffit pas~: il faut un rythme, une fa\c{c}on de stocker l'information
et un m\'ecanisme anti-abandon. Ce chapitre est court mais c'est celui qui d\'etermine si tu finis
le livre.

\section{Les cinq principes qui font la diff\'erence}

\begin{enumerate}
  \item \textbf{Pratique d\'elib\'er\'ee, pas r\'ep\'etition.} Travailler juste au-dessus de ton niveau,
        avec retour imm\'ediat (le message d'erreur, le test qui passe). Refaire un exercice
        d\'ej\`a r\'eussi est du r\'econfort, pas de l'apprentissage.
  \item \textbf{R\'ecup\'eration active.} Ferme le livre et \'ecris de m\'emoire la syntaxe d'un
        \verb|try/catch|, la diff\'erence entre \verb|==| et \verb|===|. Se tromper ancre
        dix fois mieux que relire.
  \item \textbf{Espacement.} Dix minutes de r\'evision chaque jour valent mieux que deux heures
        le dimanche. Un tour de flashcards le matin, l'erreur de la veille au soir.
  \item \textbf{Technique Feynman.} Explique \`a voix haute, avec des mots simples, comme si tu
        parlais \`a un enfant de douze ans. Tout ce que tu n'arrives pas \`a expliquer est une
        lacune --- retourne lire.
  \item \textbf{Production avant consommation.} Ratio cible~: 30 \% lire/regarder,
        70 \% \'ecrire du code. Si ta journ\'ee a \'et\'e 90 \% vid\'eo, tu ne progresses pas, tu te distraies.
\end{enumerate}

\section{La routine quotidienne (60 \`a 120 minutes)}

\begin{tabz}{llX}
  \headrow \thd{Dur\'ee} \& \thd{Bloc} \& \thd{Contenu} \\
  \midrule
  10 min \& R\'eveil \& 10 flashcards + relire le code \'ecrit hier, sans l'\'editeur. \\
  25 min \& Lecture active \& Un morceau de chapitre, le carnet ouvert, on pr\'edit les r\'esultats. \\
  50 min \& Code \& Recopier les exemples, les casser, refaire l'exercice du chapitre. \\
  15 min \& Journaliser \& \'Ecrire 5 lignes~: ce que j'ai compris, o\`u j'ai bloqu\'e, la solution. \\
  10 min \& Curiosit\'e \& Lire la doc officielle de ce que je viens d'utiliser (pas celle d'hier). \\
\end{tabz}

Si tu n'as que 25 minutes~: fais le bloc \og ~Code~ \fg{}. Jamais le bloc \og ~Lecture~ \fg{} seule.

\section{Le carnet d'erreurs --- ton actif le plus pr\'ecieux}

Cr\'ee \fichier{carnet.md} dans chaque projet et note~: \emph{texte exact de l'erreur},
\emph{contexte}, \emph{cause r\'eelle}, \emph{correction}, \emph{date}. En six mois tu auras un
manuel de d\'epannage personnel qui vaut plus que n'importe quel tuto. Les annexes
(\cref{ch-antisecche}) en sont une version g\'en\'erique~: tiens la tienne \`a jour.

\begin{codeconf}
## 2026-01-14 - PHP - 500 blank page
Erreur : aucune, page blanche.
Cause : une exception ignoree car display_errors = Off en production.
Fix : log_errors=On, error_log=php://stderr, et activer app_dev en local.
Lecon : ne jamais chercher a l'aveugle ; ouvrir les logs d'abord.
\end{codeconf}

\section{Plan de 12 semaines (version intensive)}

Adapte la vitesse~: doubler la dur\'ee de chaque ligne est parfaitementOK. Le plan n'est pas
une course, c'est un ordre.

\begin{tabz}{llX}
  \headrow \thd{Semaine} \& \thd{Objet} \& \thd{Livrable concret \`a la fin de la semaine} \\
  \midrule
  1 \& Outils, bases JS (\cref{ch-js-bases}) \& 30 petits scripts dans \verb|labo/js|, dont FizzBuzz, tri, palindrome. \\
  2 \& Objets, tableaux, cha\^ines, regex (\cref{ch-js-struct}) \& Un analyseur de texte (nombre de mots, mots les plus fr\'equents). \\
  3 \& S\'emantique avanc\'ee (\cref{ch-js-avance}) \& Un mini-syst\`eme d'\'ev\'enements \'ecrit \`a la main (20 lignes). \\
  4 \& Asynchrone (\cref{ch-js-async}) \& Un script qui t\'el\'echarge 20 pages web en parall\`ele avec limitation de concurrence. \\
  5 \& DOM et \'ev\'enements (\cref{ch-js-dom,ch-js-events}) \& Une page \og ~liste de courses~ \fg{} sans framework, avec filtre et tri. \\
  6 \& Fetch, formulaires, stockage (\cref{ch-js-dom,ch-js-storage}) \& La m\^eme liste, persist\'ee en LocalStorage + API externe. \\
  7 \& Node.js, modules, npm, CLI (\cref{ch-node,ch-js-outils}) \& Un outil de ligne de commande qui renomme en lot les fichiers d'un dossier. \\
  8 \& TypeScript + tests + s\'ecurit\'e (\cref{ch-js-ts,ch-js-tests,ch-js-secu}) \& Le projet de la semaine~6 typ\'e, test\'e \`a 80 \%, sans XSS. \\
  9 \& PHP bases, OOP, formulaires (\cref{ch-php-bases,ch-php-oop,ch-php-http}) \& Un mini-CRUD sur fichier JSON, puis sur MySQL. \\
  10 \& PDO, sessions, auth, s\'ecurit\'e PHP (\cref{ch-php-db,ch-php-http,ch-php-secu}) \& Un blog complet avec inscriptions, r\^oles, upload d'images s\'ecuris\'e. \\
  11 \& Java bases, OOP, collections, exceptions (\cref{ch-java-bases,ch-java-oop,ch-java-collections,ch-java-collections}) \& Un gestionnaire de biblioth\`eque console avec tests JUnit. \\
  12 \& Streams, JDBC, Spring Boot, d\'eploiement (\cref{ch-java-fonctionnel,ch-java-db,ch-java-spring,ch-deploy}) \& Une API REST Spring Boot connect\'ee \`a PostgreSQL, d\'eploy\'ee sur un VPS. \\
\end{tabz}

\section{Utiliser l'IA sans s'handicaper}

Un assistant de code est un excellent tuteur et un tr\`es mauvais substitut. R\`egle simple~:
\emph{l'IA r\'epond \`a tes questions, elle ne fait pas tes exercices.}

\begin{itemize}
  \item \textbf{Demander}~: \og ~explique-moi ligne par ligne ce message d'erreur~ \fg{}, \og ~quels sont les
        cas limites de cette fonction~ \fg{}, \og ~propose 3 variantes de style, avec les avantages~ \fg{}.
  \item \textbf{Refuser}~: coller l'\'enonc\'e d'un exercice et prendre la r\'eponse. Tu \'economises
        trois minutes et tu perds la comp\'etence.
  \item \textbf{Toujours}~: relire ligne par ligne le code g\'en\'er\'e et le faire tourner. Ne
        mets jamais dans un projet du code que tu ne sais pas expliquer.
  \item Outils~: \link{Copilot}{https://github.com/features/copilot},
        \link{Cursor}{https://www.cursor.com/}, \link{Claude}{https://claude.ai/},
        \link{ChatGPT}{https://chat.openai.com/}, \link{aider (CLI)}{https://aider.chat/}.
\end{itemize}

\section{Quand tu bloques --- la proc\'edure en cinq temps}

\begin{enumerate}
  \item \textbf{Lire l'erreur en entier}, y compris la ligne \og ~caused by~ \fg{}, et chercher
        le texte exact entre guillemets.
  \item \textbf{R\'eduire.} Fais un exemple minimal qui reproduit le probl\`eme (deux fichiers,
        dix lignes). Neuf fois sur dix, la r\'eduction r\'ev\`ele la cause.
  \item \textbf{Isoler une hypoth\`ese}, la tester, l'\'ecrire. \og ~Je pense que X. Je teste Y.~ \fg{}
  \item \textbf{Chercher dans la doc officielle} (les liens de fin de chapitre) plut\^ot que
        dans un blog de 2014.
  \item \textbf{Demander}, avec l'exemple minimal, la version de ton outil et le message complet
        --- sur \link{Stack~Overflow}{https://stackoverflow.com/questions/tagged/javascript},
        \link{Open~Classrooms}{https://openclassrooms.com/forum},
        \link{Dev~Discord/Slack}{https://github.com/websites/DevChatter} ou le Discord officiel
        d'un projet. Une question bien pos\'ee re\c{c}oit une r\'eponse en dix minutes.
\end{enumerate}

\begin{note}
\og ~Je tourne en rond depuis 45 minutes~ \fg{} est le signal pour lever les yeux de l'\'ecran~: marche,
douche, sommeil. Ce n'est pas une perte de temps, c'est la phase d'incubation du d\'ebogage.
\end{note}

\section{Sept familles d'erreurs que tu rencontreras}

\begin{enumerate}
  \item \textbf{Chemins}~: fichier trouv\'e sur ta machine, introuvable en production.
  \item \textbf{Encodage}~: accents cass\'es (UTF-8 partout, \verb|utf8mb4| en MySQL).
  \item \textbf{Ordre d'ex\'ecution}~: async non attendu, variable utilis\'ee avant initialisation.
  \item \textbf{Environnement}~: version d'outil, variable manquante, cache du navigateur.
  \item \textbf{Confusion de type}~: \verb|"1"+1|, \verb|0 == ""|, \verb|null| contre \verb|undefined|.
  \item \textbf{\'Etat partag\'e}~: mutation d'un objet pass\'e par r\'ef\'erence, effet de bord cach\'e.
  \item \textbf{Pr\'ematurit\'e}~: optimiser, abstraire, g\'en\'eraliser avant que le besoin existe.
\end{enumerate}

Tu peux les scanner une par une~: elles couvrent environ 80 \% des blocages de d\'ebutant.

\begin{exercice}
\begin{enumerate}
  \item \'Ecris ton contrat d'apprentissage~: jours, horaires, livrables, sanction en cas de
        manquement. Signe-le, date-le, colle-le au mur.
  \item Pr\'epare les trois fichiers \fichier{carnet.md} (un par langage) et le d\'ep\^ot
        \verb|apprentissage|.
  \item Fais une fiche de 5 cartes m\'emoire sur ce que tu as appris cette semaine
        (\link{Anki}{https://apps.ankiweb.net/} est l'outil standard, gratuit partout sauf iOS).
  \item Explique \`a voix haute, en deux minutes, la diff\'erence entre un langage interpr\'et\'e
        et un langage compil\'e. Enregistre-toi si tu oses~; r\'e\'ecoute.
\end{enumerate}
\end{exercice}

\begin{ressources}
\begin{itemize}
  \item \emph{Eloquent JavaScript} (livre gratuit de r\'ef\'erence) --- \url{https://eloquentjavascript.net/}
  \item \emph{The Road to Learn React / PHP / Java} --- \url{https://roadmap.sh/}
  \item MDN en fran\c{c}ais --- \url{https://developer.mozilla.org/fr/docs/Web/JavaScript}
  \item Exercices~: \url{https://exercism.org/}, \url{https://adventofcode.com/}, \url{https://www.codewars.com/}
  \item M\'emorisation espac\'ee~: \url{https://apps.ankiweb.net/}
\end{itemize}
\end{ressources}

\index{methode}\index{plan d'etudes}
\end{document}
